\section{LLM 的限制与适应：数据、记忆、更新和可解释性}

应用越广，模型缺口越明显。本章沿课堂顺序讨论五组问题：训练成本与数据效率，记忆和个性化，合成数据与 memorization，持续学习与 model editing，self-reflection 与 hallucination。它们都在回答同一个问题：冻结参数之后，系统怎样继续学习、更新和自我纠错。

\subsection{数据、计算与研究集中化}

训练大型模型需要海量 token、accelerator time、energy 与工程基础设施。课堂把这一点与研究民主化联系起来：当可重复实验只能由少数机构承担时，方法比较、开放审计和学生参与都会受限。本节先看成本页，再转向 BabyLM 与小模型，回答“能力提升是否只能购买更多计算”。

\lecturefigure{slide-053.jpg}{LLM 训练对数据、计算、资金与能源的需求}{官方 CS25 V4 slide deck，第 53 页}

\begin{importantbox}{训练计算的粗略尺度}
Dense Transformer 预训练常用近似
\[
C\approx 6ND,
\]
其中 $N$ 是参数量，$D$ 是训练 token 数，$C$ 是浮点运算量级。常数会随 architecture 与实现变化，但近似揭示两个事实：参数和数据都会线性推高训练计算；更高质量 token、稀疏激活与高效 kernel 可以改变实际成本。
\end{importantbox}

\begin{warningbox}{成本不只是一张 GPU 账单}
数据收集和过滤、失败实验、checkpoint、网络通信、推理服务、评估与人类反馈都消耗资源。只报最终训练 FLOPs 会低估完整研发链路，也无法说明单位能力的成本。
\end{warningbox}

\subsection{BabyLM：人类数据效率提出了什么问题}

BabyLM 页把儿童与 LLM 对比：儿童接触的语言量小得多，却能依靠感知、互动、社会反馈和具身经验形成稳健概念。这个类比不能证明机器必须复制儿童学习，但它暴露了纯文本规模化的数据效率问题。读两页对照时要关注 supervision density 和 interaction，而不是只比较词数。

\lecturefigure{slide-054.jpg}{儿童学习与 LLM 训练的数据环境差异}{官方 CS25 V4 slide deck，第 54 页}

\lecturefigure{slide-055.jpg}{BabyLM 讨论中的数据量、发展过程与多模态经验}{官方 CS25 V4 slide deck，第 55 页}

\begin{knowledgebox}{读图：儿童并不是“小数据语言模型”}
儿童拥有视觉、动作、共同注意、情绪反馈、因果干预和持续环境。把词数直接与 LLM token 比较，会忽略 supervision density 与 interaction。BabyLM 的价值在于逼迫研究者探索更好的 data curriculum、architecture 与 inductive bias。
\end{knowledgebox}

\teachervoice{课堂没有把 BabyLM 当成“儿童证明 scaling 错了”，而是把它作为研究问题：如果人类可以用更少语言数据学习，模型是否也能借助更高质量数据、不同结构或交互获得更高 sample efficiency。}

\subsection{小模型与 on-device model}

Minified LLM 与 on-device model 把目标从“最大通用能力”改成“可接受能力下的成本、延迟与隐私”。较小模型可以通过蒸馏、quantization、pruning、专用数据和 retrieval 提高实用性。

\lecturefigure{slide-056.jpg}{小型化与端侧 LLM 的动机和路线}{官方 CS25 V4 slide deck，第 56 页}

\begin{knowledgebox}{术语表：四种模型压缩路线}
\begin{tabularx}{\textwidth}{L{0.20\textwidth} X}
\toprule
方法 & 核心机制 \\
\midrule
Distillation & 用 teacher 的输出、logit 或 reasoning trace 训练 smaller student。 \\
Quantization & 用更低 bit-width 表示 weight/activation，降低内存与计算。 \\
Pruning & 删除低贡献连接、channel、head 或 block。 \\
Specialization & 用领域数据、adapter、tool 或 retrieval 缩小必须参数化的能力范围。 \\
\bottomrule
\end{tabularx}
\end{knowledgebox}

\subsection{Memory augmentation 与 personalization}

冻结 LLM 的参数不会随每次对话自动更新。Memory augmentation（记忆增强）把部分信息放到外部 store，需要时检索回 context；personalization（个性化）则让系统持续利用用户偏好、历史与约束。接下来的两页把“改输入、改外部记忆、改参数”并列，读图时先确认每条路线改变的对象。

\lecturefigure{slide-057.jpg}{冻结知识与长期个性化需求之间的矛盾}{官方 CS25 V4 slide deck，第 57 页}

\lecturefigure{slide-058.jpg}{Fine-tuning、prompt、memory bank 与 RAG 的适应路线}{官方 CS25 V4 slide deck，第 58 页}

\begin{importantbox}{四种适应机制不要互相替代}
\begin{tabularx}{\textwidth}{L{0.19\textwidth} L{0.24\textwidth} X}
\toprule
机制 & 改变什么 & 适合的问题 \\
\midrule
Prompt/context & 本次输入 & 临时任务说明与少量证据。 \\
Retrieval/RAG & 外部可更新 memory & 事实、文档、用户历史与可引用证据。 \\
Adapter/fine-tuning & 部分或全部参数 & 稳定行为、格式、领域分布与技能适应。 \\
Model editing & 定位后的局部参数 & 少量事实或关联的定向修改，需验证副作用。 \\
\bottomrule
\end{tabularx}
\end{importantbox}

\begin{importantbox}{最简 retrieval-augmented generation}
给定 query $q$、文档库 $\mathcal{D}$、检索器 $R$ 和生成器 $G$：
\[
z_{1:k}=R(q,\mathcal{D}),\qquad y=G(q,z_{1:k}).
\]
$z_{1:k}$ 是 top-$k$ evidence。检索扩大可访问知识，却不保证 evidence 正确、排序合理或 generator 真正使用它。
\end{importantbox}

\begin{lstlisting}[caption={外部记忆检索与生成的最小流程}]
def answer(query, memory, encoder, generator, top_k=5):
    query_vector = encoder(query)
    evidence = memory.search(query_vector, top_k=top_k)
    prompt = build_prompt(query=query, evidence=evidence)
    return generator(prompt), evidence
\end{lstlisting}

\teachervoice{讲者特别区分了外部 memory 与内部 skill：RAG 适合补事实和用户历史，但不会自动让模型获得新的基础能力。若 datastore 质量差，检索只会把错误更快送进 context。}

\subsection{Synthetic data 与 Phi：数据质量能否替代部分规模}

合成数据（synthetic data）用强模型生成 instruction、answer 或 reasoning trace，再训练其他模型。它降低人工收集成本，也可能复制 teacher 的错误、风格偏差和盲区。本节用 Phi 作为数据质量路线的代表案例，先看数据生成与筛选，再看小模型结果的适用范围。

\lecturefigure{slide-059.jpg}{合成预训练数据、筛选与 distillation}{官方 CS25 V4 slide deck，第 59 页}

\begin{warningbox}{Synthetic data 的三个闭环风险}
\begin{enumerate}
  \item \textbf{Error inheritance}：teacher 的错误被 student 重复学习；
  \item \textbf{Diversity collapse}：生成分布比真实世界更窄，反复训练后模式趋同；
  \item \textbf{Evaluation leakage}：prompt 或答案与 benchmark 重合，造成虚高分数。
\end{enumerate}
因此需要 provenance、deduplication、quality filter 与真实数据对照。
\end{warningbox}

\lecturefigure{slide-060.jpg}{Phi-2 用高质量与合成数据训练小模型的课堂结果}{官方 CS25 V4 slide deck，第 60 页}

\teachervoice{讲者把 Phi 的关键 takeaway 放在 data source 与 quality，而不是把小参数量本身写成魔法。更少参数能否保持能力，取决于任务覆盖、数据筛选和评估范围。}

\begin{knowledgebox}{读图：Phi 代表“每个 token 更有价值”的路线}
课堂强调 textbook-quality data 与 synthetic reasoning data，使较小模型在部分任务上接近更大模型。图不证明小模型在所有能力上等价，而是说明参数规模之外，data selection 和 curriculum 可以显著改变结果。
\end{knowledgebox}

\subsection{Learning 还是 memorizing}

训练数据与 benchmark 不透明时，模型输出可能来自组合式泛化，也可能来自近似记忆。测试污染（test contamination）会让模型在评估前见过题目或高度相似内容，使分数无法代表新任务泛化。下面的争论页应作为 evaluation design 问题阅读，而不是要求在“会思考”和“只背诵”之间仓促二选一。

\lecturefigure{slide-061.jpg}{LLM 的新知识、模式组合与 memorization 争论}{官方 CS25 V4 slide deck，第 61 页}

\begin{importantbox}{Memorization audit 的三个层次}
\begin{enumerate}
  \item \textbf{Exact match}：能否复现训练文本的长片段；
  \item \textbf{Near duplicate}：训练集是否含题目、模板或答案的轻微改写；
  \item \textbf{Behavioral generalization}：换实体、结构或约束后，机制是否仍然成立。
\end{enumerate}
只做第一层会漏掉大量 contamination，只看最终准确率又无法区分记忆与泛化。
\end{importantbox}

\begin{warningbox}{版权争议与认知争议不要混为一谈}
模型是否“理解”是认知与行为问题；训练数据是否获授权、输出是否过度复现原作是法律与治理问题。两者相关，但不能用一个结论替代另一个。
\end{warningbox}

\subsection{Continual learning：不是定期重训的别名}

Continual learning（持续学习）要求模型在新经验到来时更新，同时保留旧能力并控制灾难性遗忘。课堂用人类日常学习作对照，批评把 trace distillation 或周期 fine-tuning 直接称为永久自我改进。本节先定义 stability--plasticity tradeoff，再解释为什么单次新任务提分不足以证明持续学习。

\lecturefigure{slide-062.jpg}{持续学习、在线适应与灾难性遗忘}{官方 CS25 V4 slide deck，第 62 页}

\begin{importantbox}{Continual-learning objective}
对任务流 $\mathcal{T}_1,\ldots,\mathcal{T}_m$，更新后模型既要降低新任务损失，也要限制旧任务退化：
\[
\min_\theta\ \mathcal{L}_{\mathcal{T}_m}(\theta)+
\lambda\sum_{i<m}\Omega_i(\theta,\theta_{i}^{*}).
\]
$\Omega_i$ 是保持旧知识的 regularization、replay 或 parameter constraint，$\lambda$ 控制 plasticity 与 stability 的平衡。真正困难的是没有完整旧数据时仍能验证遗忘。
\end{importantbox}

\teachervoice{讲者用一个直观反差说明问题：人类不会每两个月坐下来重新“读完整个互联网”才能继续学习。把新 trace 蒸馏进模型更接近再训练，不等于在线、终身的持续学习。}

\subsection{Interpretability 与 model editing}

Interpretability（可解释性）试图理解内部表示和计算如何导致输出；mechanistic interpretability 更关注具体 component、circuit 与 causal intervention。Model editing 则希望定向改变某个事实关联，而不完整重训模型。

\lecturefigure{slide-063.jpg}{大模型黑箱、控制与 mechanistic interpretability}{官方 CS25 V4 slide deck，第 63 页}

\begin{knowledgebox}{三种“解释”不要混用}
\begin{tabularx}{\textwidth}{L{0.24\textwidth} X}
\toprule
层次 & 回答的问题 \\
\midrule
Behavioral explanation & 输入变化怎样影响输出，模型在什么分布失效？ \\
Attribution/probing & 哪些 token、activation 或 feature 与输出相关？ \\
Mechanistic/causal analysis & 干预某个 component 是否稳定改变目标行为？ \\
\bottomrule
\end{tabularx}
相关性可帮助定位，不自动构成因果机制。
\end{knowledgebox}

\lecturefigure{slide-064.jpg}{ROME 与 factual model editing 的课堂示意}{官方 CS25 V4 slide deck，第 64 页}

\begin{importantbox}{Rank-one edit 的抽象形式}
若希望将某层权重 $W$ 定向更新，可写成
\[
W' = W + uv^\top,
\]
其中 $u,v$ 是由目标 key/value relation 求得的向量。Rank-one update 成本低，但验收必须同时检查 edit success、paraphrase generalization、locality 与 unrelated knowledge damage。
\end{importantbox}

\begin{warningbox}{改对一条事实不代表模型“知识库已修复”}
事实可能分布在多个层和上下文中；局部编辑还可能改变相邻实体或语言模式。Model editing 需要 regression suite，而不是只展示目标 prompt 已输出新答案。
\end{warningbox}

\subsection{MoE、modularity 与自我反思}

课程再次回到 MoE，并提出一个更大胆的问题：能否像人脑功能分区一样，让单一模型内部形成可复用模块。这个类比适合激发 architecture 设计，却不能掩盖 router、expert specialization 与 global coordination 的实际难题。

\lecturefigure{slide-065.jpg}{MoE 的 routed expert 结构与稀疏激活}{官方 CS25 V4 slide deck，第 65 页}

\lecturefigure{slide-066.jpg}{从多个 expert 到单模型模块化的开放设想}{官方 CS25 V4 slide deck，第 66 页}

\begin{warningbox}{模块化必须有接口和验收指标}
如果模块之间没有明确输入输出、routing criterion、capacity 和 failure isolation，“模块化”只是一张概念图。可测指标包括 expert load、token drop、routing entropy、specialization、通信成本与故障传播。
\end{warningbox}

Self-reflection（自我反思）让模型读取自己的答案、生成 critique，再迭代修改。它可以增加 test-time compute，也能把错误暴露出来；如果 critique model 与原模型共享盲区，循环可能只产生更长、更自信的错误。

\lecturefigure{slide-067.jpg}{Self-reflection 通过 critique 与 revision 迭代输出}{官方 CS25 V4 slide deck，第 67 页}

\lecturefigure{slide-068.jpg}{多层 self-reflection 的性能变化与边界}{官方 CS25 V4 slide deck，第 68 页}

\begin{importantbox}{Reflection loop}
设初始回答为 $y_0=G(x)$，第 $t$ 轮生成 critique $c_t=C(x,y_t)$，再修订
\[
y_{t+1}=R(x,y_t,c_t).
\]
$G$ 是生成器，$C$ 是 critique 过程，$R$ 是 revision 过程。循环需要外部 stopping rule 与 verifier；否则“继续反思”可能增加成本而不增加正确性。
\end{importantbox}

\subsection{Hallucination、calibration 与外部证据}

Hallucination 页把问题描述为“模型不知道自己不知道”。更精确地说，模型可能在缺乏证据时仍给出高置信语言。Calibration（校准）关注置信度与实际正确率是否匹配；retrieval 和 verifier 则尝试引入外部证据。

\lecturefigure{slide-069.jpg}{Hallucination、事实验证、calibration 与 retrieval 路线}{官方 CS25 V4 slide deck，第 69 页}

\begin{importantbox}{Calibration 的定义}
若模型对一组回答给出置信度 $q$，理想校准要求
\[
P(\text{correct}\mid \text{confidence}=q)\approx q.
\]
校准良好不等于准确率高；一个经常回答“不确定”的模型可能很校准。反过来，高准确率模型也可能对少数错误过度自信。
\end{importantbox}

\begin{warningbox}{Temperature 不是事实性旋钮}
降低 sampling temperature 会让分布更尖、更接近 greedy decoding，却不能把错误知识变成正确知识。事实性需要 evidence、retrieval、verification、训练数据与任务约束共同支持。
\end{warningbox}

\subsection{本章小结}

LLM 的主要缺口不是一个独立 bug，而是一组相互耦合的问题：数据和计算决定可学范围，外部 memory 提供更新通道，持续学习与 model editing 改变参数，interpretability 帮助定位机制，reflection 和 calibration 改变推理时行为。任何单一方法都不能自动同时解决知识、技能、事实性和长期适应。

